\label{chap:83-21}

\input{ampliaciones_sucesoras_20260824/parte_i_ii/owners/04ac_geometria_proyectiva_bloch_rev2}

\section{Construction de l'espace de Hilbert à partir de la connexion nonadique et de ses opérateurs de phase}

\subsection{Réalisation matricielle de chaque profondeur}

Pour \(d\ge1\), nous définissons
\[
  \mathcal C_9^{(d)}=(\mathbb Z/9\mathbb Z)^d,
  \qquad
  H_d=\ell^2(\mathcal C_9^{(d)}),
  \qquad
  A_d=B(H_d)\cong M_{9^d}(\mathbb C).
\]
La base canonique \(\{\delta_x:x\in\mathcal C_9^{(d)}\}\) conserve
l’étiquette discrète de chaque état. La mesure uniforme induit la trace
normalisée
\[
  \tau_d(a)=\frac{1}{9^d}\Tr(a).
\]

La factorisation
\[
  \mathcal C_9^{(d+1)}
  =
  \mathcal C_9^{(d)}\times\mathbb Z/9\mathbb Z
\]
produit une identification unitaire
\[
  H_{d+1}\cong H_d\otimes\mathbb C^9.
\]
C’est l’entrée exacte de la tour opératorielle. Cela n’identifie pas pour autant
\((\mathbb Z/9)^3\) à \(\mathbb F_3^6\) comme groupes : ils partagent \(729\)
éléments, mais leurs opérations et cocycles de retenue sont distincts.

\subsection{Représentation de Weyl–Heisenberg de la phase et du déplacement sur la fibre tritique}

Soit \(\omega=e^{2\pi i/9}\). Sur
\(H_1=\ell^2(\mathbb Z/9\mathbb Z)\), nous définissons
\[
  (Xf)(r)=f(r-1),
  \qquad
  (Zf)(r)=\omega^r f(r).
\]
Alors
\[
  X^9=Z^9=I,
  \qquad
  ZX=\omega XZ.
\]
Le groupe engendré par \(X,Z,\omega I\) est le groupe de Heisenberg fini
d’ordre \(9^3=729\).

\begin{theorem}[Stone--von Neumann fini]
Toute représentation unitaire irréductible du groupe de Heisenberg fini dans
laquelle le centre agit par le caractère primitif
\(\omega I\) est unitairement équivalente à la représentation précédente.
\end{theorem}

\begin{proof}
Soit \(v\) un vecteur propre de \(Z\). Les vecteurs
\[
  v,Xv,\ldots,X^8v
\]
ont des valeurs propres distinctes pour \(Z\), puisque
\(ZX^kv=\omega^kX^kZv\). Ils sont orthogonaux et engendrent un sous-espace
invariant. Par irréductibilité, ils forment une base de toute la représentation.
Dans cette base, \(X\) est le déplacement cyclique et \(Z\) la phase diagonale.
\end{proof}

\begin{controlbox}
L’unicité échoue si le caractère central primitif n’est pas fixé ou si l’on
permet des représentations réductibles. Ce théorème fini ne doit pas être confondu
avec la version continue pour les relations canoniques de commutation.
\end{controlbox}

\subsection{Projecteurs arithmétiques et incompatibilité}

Les feuillets additif et multiplicatif de \(\APP\) définissent des sous-espaces
d’information et leurs projecteurs orthogonaux
\(P_{\Sigma,d},P_{\Pi,d}\). Les identités construites satisfont
\[
 \|[P_{\Sigma,2},P_{\Pi,2}]\|
 =\frac{\sqrt2}{3},
 \qquad
 \|[P_{\Sigma,3},P_{\Pi,3}]\|
 =\frac{\sqrt3}{4}.
\]
La non-commutativité ne se déduit pas d’une analogie quantique : c’est un théorème
matriciel fini. Sa lecture « pré-Heisenberg » est postérieure à ce calcul.

\begin{proposition}
Les projecteurs précédents n’appartiennent pas conjointement à une même
sous-algèbre booléenne de projections.
\end{proposition}

\begin{proof}
Dans une sous-algèbre abélienne, tous les projecteurs commutent. Les commutateurs
montrés sont non nuls.
\end{proof}

\subsection{Structure hilbertienne de la publication matricielle et mémoire généalogique préservée}

L’espace de Hilbert fournit la superposition linéaire, la projection orthogonale
et le calcul spectral. HMT apporte l’indice généalogique de la base et distingue
quel opérateur représente une transition, une mesure ou une ombre. La
réalisation ne remplace pas le \(\TPK\) : elle le représente.

\begin{sourcebox}
Les projecteurs APP et la paire Weyl--Heisenberg font partie de la construction
précédente. L’organisation \(H_d,A_d\) et la preuve d’unicité sont établies
ici par l’application du théorème classique correspondant aux objets
HMT.
\end{sourcebox}


\section{Les \(9\) phases du TPK et la géométrie continue}

\subsection{Séparation typée entre le relèvement UHF des rangs nonadiques et la prolongation géométrique de la branche marquée}

Soit \(p_j=|e_j\rangle\langle e_j|\) le projecteur sur une porte marquée.
La factorisation \(H_{d+1}\cong H_d\otimes\mathbb C^9\) admet deux
relèvements :
\[
  \iota_d(a)=a\otimes I_9,
  \qquad
  j_{d,j}(a)=a\otimes p_j.
\]

\begin{proposition}[Unité et trace]
Pour tout \(a\in A_d\),
\[
  \tau_{d+1}(\iota_d(a))=\tau_d(a),
  \qquad
  \iota_d(I)=I,
\]
tandis que
\[
  \tau_{d+1}(j_{d,j}(a))=\frac1{9}\tau_d(a),
  \qquad
  j_{d,j}(I)\ne I.
\]
\end{proposition}

\begin{proof}
Comme \(\Tr(I_9)=9\) et
\(\Tr(p_j)=1\),
\[
 \tau_{d+1}(a\otimes I_9)
 =\frac{\Tr(a)\,9}{9^{d+1}}
 =\tau_d(a),
\]
et
\[
 \tau_{d+1}(a\otimes p_j)
 =\frac{\Tr(a)}{9^{d+1}}
 =\frac{\tau_d(a)}9.
\]
Les identités d'unité sont immédiates.
\end{proof}

\begin{intuitionbox}
« Regarder les neuf phases du TPK » signifie conserver l'opérateur identité sur la
fibre complète. Une porte isolée conserve une histoire marquée ; les neuf
ensemble conservent en outre l'unité et la normalisation globale.
\end{intuitionbox}

\subsection{La limite UHF}

\begin{theorem}[Clôture UHF nonadique]
La limite inductive unitale
\[
  A_{9^\infty}
  :=
  \varinjlim(A_d,\iota_d)
  \cong
  \bigotimes_{n=1}^{\infty}M_9(\mathbb C)
\]
est l'algèbre UHF de type surnaturel \(9^\infty=3^\infty\). Elle possède une
trace produit \(\tau\) compatible avec toutes les \(\tau_d\).
\end{theorem}

\begin{proof}
Les inclusions sont injectives, unitales et préservent la trace. La réunion
algébrique est formée de matrices finies emboîtées ; sa fermeture en norme est,
par définition, une UHF. La trace produit est l'unique trace compatible avec les
traces matricielles de tous les facteurs \citep{glimm1960}.
\end{proof}

\begin{theorem}[Facteur hyperfini]
Soit \((\pi_\tau,\mathcal H_\tau,\Omega_\tau)\) la représentation GNS de la
trace. Alors
\[
  \pi_\tau(A_{9^\infty})''\cong\Rhyper,
\]
le facteur hyperfini de type \(II_1\).
\end{theorem}

\begin{proof}
La fermeture est finie puisque \(\tau(I)=1\). L'extrémalité de la trace en fait un
facteur. Il est hyperfini puisque la réunion croissante des matrices
\(\pi_\tau(A_d)\) est faiblement dense. L'unicité du facteur séparable
hyperfini \(II_1\) complète l'identification
\citep{murrayvonneumann1936,connes1976}.
\end{proof}

\subsection{La branche marquée change de type}

Soit \(V_{d,j}:H_d\to H_{d+1}\), \(V_{d,j}\xi=\xi\otimes e_j\).
Alors
\[
  j_{d,j}(a)=V_{d,j}aV_{d,j}^*.
\]
À la limite des espaces, la réunion des coins contient les
opérateurs de rang fini. Sa fermeture en norme est
\(\mathcal K(H_\infty)\), et le bicommutant naturel est
\(B(H_\infty)\), de type \(I_\infty\).

\begin{figure}[H]
\centering
\resizebox{\textwidth}{!}{%
\begin{tikzpicture}[>=Latex,node distance=12mm and 17mm,
  box/.style={rounded corners,draw=hmtblue,fill=hmtlightblue,
    minimum width=34mm,minimum height=13mm,align=center},
  gold/.style={rounded corners,draw=hmtgold,fill=hmtlightgold,
    minimum width=34mm,minimum height=13mm,align=center}]
  \node[box] (a) {\(M_{9^d}\)\\\scriptsize état fini};
  \node[box,above right=of a] (full) {\(a\otimes I_9\)\\\scriptsize neuf phases du TPK};
  \node[gold,below right=of a] (one) {\(a\otimes p_j\)\\\scriptsize porte marquée};
  \node[box,right=of full] (uhf) {\(\mathrm{UHF}(9^\infty)\)\\\scriptsize unital/tracial};
  \node[box,right=of uhf] (ii) {\(\Rhyper\)\\\scriptsize type \(II_1\)};
  \node[gold,right=of one] (k) {\(\mathcal K(H_\infty)\)\\\scriptsize compacts};
  \node[gold,right=of k] (i) {\(B(H_\infty)\)\\\scriptsize type \(I_\infty\)};
  \draw[->,very thick,hmtblue] (a)--(full);
  \draw[->,very thick,hmtblue] (full)--(uhf);
  \draw[->,very thick,hmtblue] (uhf)--(ii);
  \draw[->,very thick,hmtgold] (a)--(one);
  \draw[->,very thick,hmtgold] (one)--(k);
  \draw[->,very thick,hmtgold] (k)--(i);
  \node[above=2mm of full,text=hmtgreen]{\scriptsize unité et trace};
  \node[below=2mm of one,text=hmtred]{\scriptsize trace \(\div9\)};
\end{tikzpicture}
}
\caption{Une porte marquée et les neuf phases du TPK produisent des fermetures de types
distincts.}
\label{p12-83-c21-s2:fig:operator-gates}
\end{figure}

\subsection{Dimension continue à partir de rangs nonadiques}

Pour une projection \(P\in M_{9^d}\),
\[
  d_d(P)=\tau_d(P)=\frac{\rank(P)}{9^d}.
\]
Les valeurs appartiennent à
\(\mathbb Z[1/9]\cap[0,1]\), ensemble dense dans \([0,1]\).

\begin{theorem}[Géométrie continue discrète]
Pour chaque \(t\in[0,1]\), il existe une projection \(P_t\in\Rhyper\) avec
\(\tau(P_t)=t\). En outre,
\[
  P\sim Q
  \quad\Longleftrightarrow\quad
  \tau(P)=\tau(Q),
\]
où \(\sim\) est l'équivalence de Murray--von Neumann.
\end{theorem}

\begin{proof}
Soit \(m_d=\lfloor9^dt\rfloor\). Puisque
\(0\le m_{d+1}-9m_d\le8\), on peut choisir des projections emboîtées de rang
\(m_d\). Leur limite forte \(P_t\) satisfait, par normalité de la trace,
\[
  \tau(P_t)=\lim_d\frac{m_d}{9^d}=t.
\]
La classification des projections dans un facteur fini donne l'équivalence
finale.
\end{proof}

La dimension continue n'est pas une cardinalité fractionnaire. Elle mesure l'équivalence des
sous-espaces dans la fermeture. Chaque approximant conserve un rang entier ; la
continuité apparaît lorsque l'on ferme toute la généalogie unitale.

\begin{resultbox}
La conservation holographique de l'unité devient un critère
opératoriel réfutable : remplacer \(I_9\) par \(p_j\) divise la trace et change
le type de la limite. L'ensemble des portes n'est pas un recensement ; c'est la
condition de l'inclusion unitale et traciale.
\end{resultbox}


\section{Système dynamique profini : translation, algèbre logique et espace des états}

\subsection{La 2ᵉ route vers le facteur hyperfini}

La construction HMT produit le groupe compact
\[
  K_\infty\cong C_4\times\mathbb Z_3
\]
comme limite de noyaux de norme de Pell, avec une translation
\(T_u(x)=x+u\) minimale et uniquement ergodique pour la mesure de Haar
\(\mu_H\). Nous considérons
\[
  M_{\rm clk}
  =
  L^\infty(K_\infty,\mu_H)\rtimes_{T_u}\mathbb Z.
\]

\begin{theorem}[Facteur de l'horloge]
Si \(T_u\) est libre, ergodique et préserve Haar, alors
\[
  M_{\rm clk}\cong\Rhyper.
\]
\end{theorem}

\begin{proof}
L'action libre et ergodique p.m.p. produit un facteur fini de trace
\[
 \tau_{\rm clk}\!\left(\sum_nf_nV^n\right)
 =\int f_0\,d\mu_H.
\]
La moyennabilité de \(\mathbb Z\) rend hyperfinie la relation d'orbites ;
le facteur est donc le facteur hyperfini \(II_1\).
\end{proof}

\begin{figure}[H]
\centering
\resizebox{\textwidth}{!}{%
\begin{tikzpicture}[>=Latex,node distance=12mm and 15mm,
  box/.style={rounded corners,draw=hmtblue,fill=hmtlightblue,
    minimum width=36mm,minimum height=13mm,align=center},
  boundary/.style={rounded corners,draw=hmtred,dashed,fill=white,
    minimum width=36mm,minimum height=13mm,align=center}]
  \node[box] (tower) {\((\mathbb Z/9)^d\)\\\(M_{9^d}\)};
  \node[box,right=of tower] (uhf) {\(\mathrm{UHF}(9^\infty)\)};
  \node[box,right=of uhf] (r1) {\(\Rhyper\)};
  \node[box,below=18mm of tower] (clock) {\(C_4\times\mathbb Z_3\)\\\(T_u\)};
  \node[box,right=of clock] (cross) {\(L^\infty(K_\infty)\rtimes\mathbb Z\)};
  \node[box,right=of cross] (r2) {\(\Rhyper\)};
  \node[boundary,right=16mm of r1,yshift=-15mm] (theta)
    {\(\Theta\) HMT canonique\\\scriptsize mesure, retenue, phase};
  \draw[->,thick,hmtblue] (tower)--(uhf);
  \draw[->,thick,hmtblue] (uhf)--(r1);
  \draw[->,thick,hmtblue] (clock)--(cross);
  \draw[->,thick,hmtblue] (cross)--(r2);
  \draw[<->,very thick,hmtgreen] (r1)--node[right]{\scriptsize isomorphie abstraite}(r2);
  \draw[dashed,->,hmtred] (r1)--(theta);
  \draw[dashed,->,hmtred] (r2)--(theta);
\end{tikzpicture}
}
\caption{Deux routes abstraites vers \(\Rhyper\). La conjugaison causale HMT
est un problème supplémentaire.}
\label{p12-83-c21-s3:fig:two-routes-r}
\end{figure}

La coïncidence des fermetures faibles n'identifie pas les antécédents
\(C^*\). La branche UHF et le produit croisé de l'odomètre peuvent posséder
des invariants \(K\)-théoriques distincts. Une identification HMT canonique exige
une application \(\Theta\) avec
\[
  \Theta_*\mu_{\APP}=\mu_H,
  \qquad
  \Theta S_{\rm carry}=T_u\Theta,
\]
et la préservation des diagonales, de la phase, de l'orientation et des projecteurs. Cet
entrelaceur demeure une frontière.

\subsection{Théorème ergodique moyen de von Neumann pour l'odomètre de Haar dans \(K_\infty\)}

Soit \(U f=f\circ T_u^{-1}\) dans \(L^2(K_\infty,\mu_H)\). Le théorème
ergodique moyen de von Neumann \citep{vonneumann1932ergodic} donne
\[
 \frac1N\sum_{j=0}^{N-1}U^jf
 \longrightarrow
 P_{\mathrm{Fix}(U)}f.
\]
L'ergodicité identifie la limite à
\[
  \left(\int f\,d\mu_H\right)\mathbf1.
\]
L'horloge possède un spectre purement ponctuel et une entropie nulle, mais ses moyennes
temporelles convergent vers la moyenne de Haar. Aucun mélange n'est nécessaire.

\subsection{Treillis orthomodulaire des projections et logique des sous-espaces fermés}

Le treillis \(\Proj(\Rhyper)\), avec
\[
  P^\perp=I-P,
  \qquad P\wedge Q,
  \qquad P\vee Q=(P^\perp\wedge Q^\perp)^\perp,
\]
est complet et orthomodulaire. Il n'est pas distributif. Les projecteurs
non commutatifs d'APP s'insèrent dans \(\Rhyper\) et fournissent un témoin
fini de cette non-distributivité \citep{birkhoffvonneumann1936}.

Chaque sous-algèbre abélienne sélectionne un contexte booléen. La logique classique
ne disparaît pas : elle occupe les contextes commutatifs. Ce qui échoue globalement est
la possibilité de réunir toutes les observables incompatibles dans une seule
algèbre booléenne.

\subsection{Matrices de densité et entropie}

Dans \(M_{9^d}\), un état normal s'écrit
\[
  \varphi_\rho(a)=\Tr(\rho a),
  \qquad \rho\ge0,\quad\Tr\rho=1.
\]
Dans \(\Rhyper\), une densité normale est
\[
  h\in L^1(\Rhyper,\tau)_+,
  \qquad \tau(h)=1,
\]
et \(\varphi_h(a)=\tau(ha)\). Son entropie relative à la trace est
\[
  S_\tau(h)=-\tau(h\log h).
\]

Si \(h=9^d\rho\), alors
\[
  S_\tau(h)=S_{\rm vN}(\rho)-d\log9.
\]
L'extension canonique
\[
  \rho_{d+1}=\rho_d\otimes\frac{I_9}{9}
\]
ajoute \(\log9\) à l'entropie matricielle, mais laisse invariant
\(S_\tau(h)\). La première quantité enregistre une fibre uniforme supplémentaire ;
la seconde utilise la normalisation holographique de la tour.

\begin{controlbox}
On n'utilisera pas indifféremment \(\Tr\) et \(\tau\). On n'interprétera pas non plus
\(\tau(P)\) comme un cardinal de points, et l'on n'appellera pas pur tout vecteur d'une
représentation. Un facteur diffus \(II_1\) ne possède pas de projections minimales.
\end{controlbox}

\subsection{Types opératoriels et frontière modulaire}

La classification construite jusqu'ici est :
\[
\begin{array}{c|c}
\text{objet}&\text{type}\\
\hline
M_{9^d}&I_{9^d}\\
\text{branche marquée}&I_\infty\\
\text{branche UHF traciale}&II_1\\
\text{horloge p.m.p.}&II_1.
\end{array}
\]
Un type \(III\) exigerait une mesure quasi invariante avec un cocycle de
Radon--Nikodym non trivial ou un état KMS dont la représentation GNS ne serait pas
semi-finie. La simple présence d'un mot « KMS » ne sélectionne pas le type.


\begin{workbox}
La complétion du groupoïde total de mémoire du \TPK\ exige une topologie
ou une structure borélienne, une mesure, la fermeture par inversion, la moyennabilité,
l'ergodicité et le contrôle de l'isotropie. Le sous-groupoïde de l'horloge satisfait ces
hypothèses ; le groupoïde total n'est pas promu par analogie.
\end{workbox}
